\documentclass[11pt,a4paper]{report}
\usepackage[utf8]{inputenc}
\usepackage[T1]{fontenc}
\usepackage{lmodern,microtype}
\usepackage[margin=2.4cm]{geometry}
\usepackage{graphicx}
\usepackage{booktabs,longtable,array,amsmath,amssymb,xcolor,listings,hyperref,enumitem,titlesec,etoolbox}
\hypersetup{colorlinks=true,linkcolor=blue!50!black,urlcolor=blue!50!black,citecolor=blue!50!black,pdftitle={nwergm: Exponential random graph models in Stata. Manual},pdfauthor={Thomas U. Grund}}
\lstdefinestyle{stata}{basicstyle=\ttfamily\footnotesize,backgroundcolor=\color{black!4},frame=single,rulecolor=\color{black!20},breaklines=true,columns=fullflexible,keepspaces=true,xleftmargin=2pt,xrightmargin=2pt,aboveskip=8pt,belowskip=8pt}
\lstset{style=stata}
\setlength{\emergencystretch}{3em}
\newcommand{\nwergm}{\texttt{nwergm}}
\newcommand{\ergm}{\texttt{ergm}}
% \opt sets an option name in typewriter type and lets long names break after ( : , and _ so that they do not run into the next table column
\ExplSyntaxOn
\tl_new:N \l__nwergm_opt_tl
\NewDocumentCommand{\opt}{m}
  {
    \tl_set:Nn \l__nwergm_opt_tl {#1}
    \regex_replace_all:nnN { ([\:\,]) } { \1 \c{allowbreak}\cB\{\cE\} } \l__nwergm_opt_tl
    \regex_replace_all:nnN { \c{\_} } { \c{\_} \c{allowbreak}\cB\{\cE\} } \l__nwergm_opt_tl
    \regex_replace_all:nnN { ([a-z])([A-Z]) } { \1 \c{allowbreak}\cB\{\cE\} \2 } \l__nwergm_opt_tl
    \regex_replace_all:nnN { (/) } { \1 \c{allowbreak}\cB\{\cE\} } \l__nwergm_opt_tl
    \regex_replace_all:nnN { ([\|\(\[=]) } { \1 \c{allowbreak}\cB\{\cE\} } \l__nwergm_opt_tl
    \texttt{ \tl_use:N \l__nwergm_opt_tl }
  }
\ExplSyntaxOff
% the listing of the output of an example (manual/ex/out/<name>.txt, written by run_examples.sh)
\newcommand{\showlog}[1]{\lstinputlisting{ex/out/#1.txt}}
\newcommand{\showlogpart}[3]{\lstinputlisting[firstline=#2,lastline=#3]{ex/out/#1.txt}}
\newcommand{\inc}[1]{\typeout{>>>#1}\input{#1}}
\newcommand{\ergmver}{4.12.0}\newcommand{\nErgm}{139}\newcommand{\nImpl}{123}\newcommand{\nPartial}{2}\newcommand{\nCovered}{125}\newcommand{\nNotPlanned}{14}\newcommand{\nValued}{13}\newcommand{\nBinary}{126}\newcommand{\pctCovered}{89.9}\newcommand{\nOther}{32}\newcommand{\nTergm}{15}\newcommand{\nMulti}{15}\newcommand{\nCount}{2}\newcommand{\nAllDoc}{171}

\newcommand{\cmpStatMax}{1.5e-08}\newcommand{\cmpMpleMax}{4.1e-08}\newcommand{\cmpSeRel}{1.0e-04}\newcommand{\cmpMcmleSd}{0.07}\newcommand{\cmpSimRel}{0.03}

\newcommand{\manualseries}{1.2}
% The manual is labelled by the minor series (valid for 1.2.x): change \manualseries only for a new minor or major release whose documented behaviour changed.
\newcommand{\pkgversion}{1.3.0}
\newcommand{\pkgdate}{26 October 2026}
% Exact version and date of the release that produced the listings (written by tools/set_version.py; rerun the examples when it changes).
% DOI of the archived manual (Zenodo). Empty until the manual is archived; nothing is printed for an empty DOI.
\newcommand{\manualdoi}{10.5281/zenodo.23661332}% UNCONFIRMED: version DOI or all-versions DOI; replace by the all-versions DOI after publication
\newcommand{\citedoi}{\ifdefempty{\manualdoi}{}{ Zenodo. \url{https://doi.org/\manualdoi}}}
\title{\textbf{nwergm}\\[4pt]\large Exponential random graph models in Stata\\ (a package that builds on \texttt{nwcommands})\\[6pt]\normalsize User manual and reference}
\author{Thomas U. Grund\\ RWTH Aachen University}
\date{Manual for version \manualseries\\ (examples produced with \texttt{nwergm} \pkgversion, \pkgdate)}
\begin{document}
\maketitle
\begin{quote}\noindent\textbf{How to cite.} This manual describes the series~\manualseries{} of \nwergm{} (it holds for all versions 1.2.x); the listings were produced with version~\pkgversion{} (\pkgdate). It is the citable reference for the package. Please cite it as:\par\smallskip
\noindent Grund, T.~U. (2026). \emph{nwergm: Exponential random graph models in Stata. Manual for version \manualseries.}\citedoi\par\smallskip
\noindent Please also report the exact version of the program that produced your results, as printed by \texttt{nwergm version} and stored in \texttt{e(version)}.
\end{quote}
\chapter*{Acknowledgments}\addcontentsline{toc}{chapter}{Acknowledgments}
\nwergm{} follows the statnet project closely and exists because of it. The exponential random graph model and the \ergm{} package are due to Mark S.\ Handcock, David R.\ Hunter, Carter T.\ Butts, Steven M.\ Goodreau, Pavel N.\ Krivitsky, Martina Morris and many contributors~\cite{hunter2008jss,handcock2008}. \ergm{}~\ergmver{} is the reference implementation on which the terms, the estimation algorithms, the defaults and the output of \nwergm{} are modeled. The R source code and the manuals of \ergm{} were the specification against which the Stata code was written and checked. The curved exponential families, the pseudolikelihood and the goodness-of-fit method of the package are those of Hunter and Handcock~\cite{hunter2006} and Hunter, Goodreau and Handcock~\cite{hunter2008jasa}; the treatment of missing dyads is that of Handcock and Gile~\cite{handcock2010}.

\nwergm{} is not affiliated with or endorsed by the statnet project, and its errors are the author's. It is meant for researchers who work in Stata or prefer it to R. It is not a replacement for \ergm{}, which remains the reference and the more complete program.

\paragraph{Use of an AI coding assistant.} The software was developed with Claude Code (Anthropic), an AI coding assistant, working under the author's direction. It helped port algorithms from ergm into Stata, Mata and C and helped write the tests. The author specified the scope, reviewed the code, and is responsible for the package. Its correctness rests on validation, not on the assistant: the results are compared numerically with ergm, and the test suite reruns these comparisons.

\tableofcontents
\inc{ch01_intro}
\inc{ch02_model}
\inc{ch03_start}
\inc{ch04_options}
\inc{ch05_terms}
\inc{ch06_special}
\inc{ch07_postest}
\inc{ch08_examples}
\inc{ch09_validation}
\inc{ch10_limits}
\appendix
\inc{app_translation}
\bibliographystyle{plain}
\begin{thebibliography}{99}
\bibitem{hunter2008jss} Hunter, D.R., Handcock, M.S., Butts, C.T., Goodreau, S.M., Morris, M. (2008). ergm: A package to fit, simulate and diagnose exponential-family models for networks. \emph{Journal of Statistical Software} 24(3). \url{https://doi.org/10.18637/jss.v024.i03}
\bibitem{handcock2008} Handcock, M.S., Hunter, D.R., Butts, C.T., Goodreau, S.M., Morris, M. (2008). statnet: Software tools for the representation, visualization, analysis and simulation of network data. \emph{Journal of Statistical Software} 24(1). \url{https://doi.org/10.18637/jss.v024.i01}
\bibitem{hunter2006} Hunter, D.R., Handcock, M.S. (2006). Inference in curved exponential family models for networks. \emph{Journal of Computational and Graphical Statistics} 15(3), 565--583. \url{https://doi.org/10.1198/106186006X133069}
\bibitem{hunter2008jasa} Hunter, D.R., Goodreau, S.M., Handcock, M.S. (2008). Goodness of fit of social network models. \emph{Journal of the American Statistical Association} 103, 248--258. \url{https://doi.org/10.1198/016214507000000446}
\bibitem{handcock2010} Handcock, M.S., Gile, K.J. (2010). Modeling social networks from sampled data. \emph{Annals of Applied Statistics} 4(1), 5--25. \url{https://doi.org/10.1214/08-AOAS221}
\bibitem{krivitsky2014} Krivitsky, P.N., Handcock, M.S. (2014). A separable model for dynamic networks. \emph{Journal of the Royal Statistical Society, Series B} 76, 29--46. \url{https://doi.org/10.1111/rssb.12014}
\bibitem{snijders2002} Snijders, T.A.B. (2002). Markov chain Monte Carlo estimation of exponential random graph models. \emph{Journal of Social Structure} 3(2), 1--40.
\bibitem{robins2007} Robins, G., Pattison, P., Kalish, Y., Lusher, D. (2007). An introduction to exponential random graph ($p^*$) models for social networks. \emph{Social Networks} 29, 173--191. \url{https://doi.org/10.1016/j.socnet.2006.08.002}
\bibitem{wasserman1996} Wasserman, S., Pattison, P. (1996). Logit models and logistic regressions for social networks: I. An introduction to Markov graphs and $p^*$. \emph{Psychometrika} 61, 401--425. \url{https://doi.org/10.1007/BF02294547}
\bibitem{frank1986} Frank, O., Strauss, D. (1986). Markov graphs. \emph{Journal of the American Statistical Association} 81, 832--842. \url{https://doi.org/10.1080/01621459.1986.10478342}
\bibitem{hummel2012} Hummel, R.M., Hunter, D.R., Handcock, M.S. (2012). Improving simulation-based algorithms for fitting ERGMs. \emph{Journal of Computational and Graphical Statistics} 21(4), 920--939. \url{https://doi.org/10.1080/10618600.2012.679224}
\bibitem{geyer1992} Geyer, C.J., Thompson, E.A. (1992). Constrained Monte Carlo maximum likelihood for dependent data. \emph{Journal of the Royal Statistical Society, Series B} 54(3), 657--683. \url{https://doi.org/10.1111/j.2517-6161.1992.tb01443.x}
\bibitem{hunter2007} Hunter, D.R. (2007). Curved exponential family models for social networks. \emph{Social Networks} 29(2), 216--230. \url{https://doi.org/10.1016/j.socnet.2006.08.005}
\bibitem{morris2008} Morris, M., Handcock, M.S., Hunter, D.R. (2008). Specification of exponential-family random graph models: terms and computational aspects. \emph{Journal of Statistical Software} 24(4). \url{https://doi.org/10.18637/jss.v024.i04}
\bibitem{krivitsky2012} Krivitsky, P.N. (2012). Exponential-family random graph models for valued networks. \emph{Electronic Journal of Statistics} 6, 1100--1128. \url{https://doi.org/10.1214/12-EJS696}
\bibitem{lazega2001} Lazega, E. (2001). \emph{The Collegial Phenomenon: The Social Mechanisms of Cooperation Among Peers in a Corporate Law Partnership}. Oxford: Oxford University Press.
\bibitem{padgett1993} Padgett, J.F., Ansell, C.K. (1993). Robust action and the rise of the Medici, 1400--1434. \emph{American Journal of Sociology} 98, 1259--1319. \url{https://doi.org/10.1086/230190}
\bibitem{kapferer1972} Kapferer, B. (1972). \emph{Strategy and Transaction in an African Factory: African Workers and Indian Management in a Zambian Town}. Manchester: Manchester University Press.
\end{thebibliography}
\end{document}
